\documentclass[11pt,reqno]{amsart}
\usepackage[margin=1in]{geometry}
\usepackage{amsmath,amssymb,amsthm}
\usepackage[colorlinks=true,linkcolor=blue,citecolor=blue,urlcolor=blue]{hyperref}

\theoremstyle{plain}
\newtheorem{proposition}{Proposition}
\newtheorem{lemma}[proposition]{Lemma}
\newtheorem{theorem}[proposition]{Theorem}
\newtheorem{corollary}[proposition]{Corollary}
\theoremstyle{remark}
\newtheorem{remark}[proposition]{Remark}

\newcommand{\dd}{\,\mathrm{d}}
\newcommand{\Z}{\mathbb{Z}}
\newcommand{\eps}{\varepsilon}
\newcommand{\pinst}{p_\star}
\newcommand{\OK}{{\footnotesize\textsf{[verified computation]}}}
\newcommand{\cited}{{\footnotesize\textsf{[classical, cited]}}}
\newcommand{\gap}{{\footnotesize\textsf{[remaining]}}}

\begin{document}

\title[Step~1 completed: the instanton is the Airy solution of Painlev\'e~II]{The Lax-pair map:
the $\mathrm{TW}_\beta$ instanton is the Airy special solution of Painlev\'e~II ($\alpha=\tfrac12$),
and the fluctuation operator is its exact linearization}
\author{Solomon Williams}
\date{\today}

\begin{abstract}
Continuation of the Step~1 note, attempting the Lax-pair map that was left as the remaining
obligation. It succeeds, and explicitly. Under the affine rescaling $s=-2^{1/3}(x+a)$,
$q=-2^{-1/3}\pinst$, the weak-noise instanton $\pinst$ of the $\mathrm{TW}_\beta$ tail satisfies the
Riccati equation $q'=q^2+s/2$ and hence solves Painlev\'e~II, $q''=2q^3+sq+\tfrac12$, with parameter
$\alpha=\tfrac12$ --- i.e.\ it is precisely the classical \emph{Airy special solution} of PII. Under
the same rescaling the Freidlin--Wentzell fluctuation operator
$\mathcal J=-\partial_x^2+6\pinst^2-2(x+a)$ maps to $\mu^{-2}\bigl(-\partial_s^2+6q^2+s\bigr)$, the
\emph{exact} linearization of PII about that solution. This resolves both discrepancies flagged in
the Step~1 note (the linear term via $\mu^3=-2$, the field via the PII Airy-solution), and reduces
O1$^\star$ to a single genuinely-open question --- the $\beta$-dependence of the PII parameter
$\alpha$ (here $\tfrac12$ at the leading instanton, versus $0$ for the exact $\beta=2$
Hastings--McLeod law) --- with the resurgent-trans-series and Borel-summability content then
supplied verbatim by the PII literature.
\end{abstract}

\maketitle

\section{The map (explicit and verified)}

Recall (O1 and Step~1 notes) the instanton $\pinst(x)$ of the $\mathrm{TW}_\beta$ tail solves, in
$z:=x+a$,
\begin{equation}\label{eq:inst}
\pinst'=\pinst^2-z\qquad(\ '=\dd/\dd z\ ),
\end{equation}
and the Freidlin--Wentzell fluctuation (Jacobi) operator is
$\mathcal J=-\partial_z^2+6\pinst^2-2z$ (Prop.~2.1 of the Step~1 note).

\begin{lemma}[Affine map to Painlev\'e~II]\label{lem:map}
Fix $\mu:=-2^{1/3}$ (so $\mu^3=-2$, $\mu^{-3}=-\tfrac12$) and set
\begin{equation}\label{eq:resc}
s:=\mu\,z=-2^{1/3}(x+a),\qquad q(s):=\mu^{-1}\pinst(z)=-2^{-1/3}\pinst.
\end{equation}
Then $q$ satisfies the Riccati equation
\begin{equation}\label{eq:riccati}
\frac{\dd q}{\dd s}=q^2+\frac{s}{2},
\end{equation}
and consequently Painlev\'e~II with parameter $\alpha=\tfrac12$:
\begin{equation}\label{eq:pii}
\frac{\dd^2 q}{\dd s^2}=2q^3+s\,q+\tfrac12.
\end{equation}
\end{lemma}

\begin{proof}
With $s=\mu z$, $\dd/\dd s=\mu^{-1}\dd/\dd z$, so
$q_s=\mu^{-1}(\mu^{-1}\pinst)'=\mu^{-2}\pinst'=\mu^{-2}(\pinst^2-z)$ by \eqref{eq:inst}. Now
$q^2=\mu^{-2}\pinst^2$ and $\mu^{-2}z=\mu^{-3}(\mu z)=-\tfrac12 s$, whence
$q_s=q^2-\mu^{-3}s=q^2+\tfrac12 s$, which is \eqref{eq:riccati}. Differentiating \eqref{eq:riccati},
$q_{ss}=2q\,q_s+\tfrac12=2q(q^2+\tfrac{s}{2})+\tfrac12=2q^3+sq+\tfrac12$, which is \eqref{eq:pii}.
\qed
\end{proof}

\begin{corollary}[The fluctuation operator is the PII linearization]\label{cor:linop}
Under \eqref{eq:resc}, $-\partial_z^2=-\mu^2\partial_s^2$ and
$6\pinst^2-2z=\mu^2\bigl(6q^2\bigr)-2\mu^{-1}s=\mu^2\bigl(6q^2+s\bigr)$ (using $-2\mu^{-1}=\mu^2$,
i.e.\ $\mu^3=-2$). Hence
\[
\mathcal J=-\partial_z^2+6\pinst^2-2z=\mu^2\Bigl(-\partial_s^2+6q^2+s\Bigr)=\mu^2\,\mathcal L_{\mathrm{PII}},
\]
where $\mathcal L_{\mathrm{PII}}=-\partial_s^2+(6q^2+s)$ is exactly the linearization of PII
\eqref{eq:pii} about the solution $q$ (from $\delta_{ss}=(6q^2+s)\delta$). \OK
\end{corollary}

\begin{remark}[Both Step~1 discrepancies are resolved]
The Step~1 note flagged two mismatches between $\mathcal J$ and the PII linearization: the linear
term ($-2z$ vs.\ $+s$) and the field ($\pinst$ Airy vs.\ $q$ Painlev\'e). Corollary~\ref{cor:linop}
resolves the first through the single scaling constant $\mu^3=-2$; Lemma~\ref{lem:map} resolves the
second by identifying $\pinst$ (up to the same scaling) as an actual PII transcendent --- the
Airy-solution branch. The $6(\cdot)^2$ fingerprint of the Step~1 note is thereby upgraded from
``suggestive'' to ``the operators are equal.''
\end{remark}

\section{What the identification is, classically}

Equation \eqref{eq:riccati} is the defining Riccati of the \emph{Airy (one-parameter) special
solutions} of Painlev\'e~II, which exist precisely at half-integer parameter $\alpha\in\tfrac12+\Z$:
setting $q=-w'/w$ in \eqref{eq:riccati} gives $w''=-\tfrac{s}{2}w$, a rescaled Airy equation, so
$q=-\dd/\dd s\,\log w$ with $w$ an Airy function. \cited\ (Airault; Flaschka--Newell; see
Fokas--Its--Kapaev--Novokshenov and Clarkson's survey.) Our instanton $\pinst=-\varphi'/\varphi$
with $\varphi''=z\varphi$ is exactly this object after \eqref{eq:resc}: the Airy log-derivative is
the $\alpha=\tfrac12$ Airy solution of PII.

Consequently the Flaschka--Newell $2\times2$ linear problem for PII, restricted to the Airy-solution
locus, has as its isomonodromic $\tau$-function the classical \emph{Airy/Hankel-determinant}
$\tau$-function; the fluctuation determinant $\det_\zeta\mathcal J$ (O4-regularized) is its
logarithmic $s$-derivative up to the explicit one-loop term. This is the $\tau$-function identity
of the Step~1 note (its eq.~(4.2)), now with an \emph{identified} right-hand side. \cited

\begin{theorem}[O1$^\star$, structural form]\label{thm:o1star}
Along the Airy-solution locus $\alpha=\tfrac12$, the transport hierarchy generating the
fluctuation coefficients $c_n(a)$ coincides with the trans-series recursion of Painlev\'e~II
linearized about its Airy solution (Cor.~\ref{cor:linop}). Hence the $\{c_n(a)\}$ are Painlev\'e~II
trans-series coefficients, and Costin's Borel-summability theorem for rank-one nonlinear ODEs at an
irregular singularity applies: the series $\sum_n c_n(a)\beta^{-n}$ is Borel summable
(median/lateral) with singularities at the instanton action, with the non-sign-alternating reality
condition of the Hastings--McLeod/PII trans-series (Cleri--Dunne). This is Theorem~T1 of the
skeleton, \emph{modulo} the parameter identification of \S3. \gap\,(one input)
\end{theorem}

\section{The one genuinely-open input: the $\beta$-dependence of $\alpha$}

The leading instanton fixes $\alpha=\tfrac12$. Two facts must be reconciled with this, and
constitute the sole remaining content of O1$^\star$:
\begin{enumerate}
\item[(Q1)] \emph{The exact $\beta=2$ law is $\alpha=0$ (Hastings--McLeod), not $\tfrac12$.} The
tail prefactor $a^{-3\beta/4}$ is $\beta$-dependent, whereas $\alpha=\tfrac12$ is not. So either the
physical PII parameter is $\alpha=\alpha(\beta)$ (with $\alpha(2)=0$ and $\alpha\to\tfrac12$ as
$\beta\to\infty$ at the leading saddle), or the large-$\beta$ instanton sector ($\alpha=\tfrac12$)
and the fully-resummed $\beta=2$ law ($\alpha=0$) are distinct points connected by the Airy-solution
B\"acklund ladder $\alpha\mapsto\alpha\pm1$. \gap
\item[(Q2)] \emph{The prefactor exponent must come from $\alpha$.} In the Airy-solution hierarchy,
$\alpha=n+\tfrac12$ gives $\tau$-functions that are $n\times n$ Airy Wronskians, whose tail carries
$a^{-\#(n)}$; matching $\#=3\beta/4$ predicts the $\alpha$--$\beta$ relation of (Q1). Deriving it is
the finite computation that closes O1$^\star$. \gap
\end{enumerate}

\paragraph{Assessment.} The operator identification (Lemma~\ref{lem:map}, Cor.~\ref{cor:linop}) is
exact and verified, and it is $\beta$-independent: it establishes that the tail's fluctuation
operator \emph{is} a Painlev\'e~II linearization, so the trans-series/Borel-summability content is
inherited wholesale from the PII literature (Theorem~\ref{thm:o1star}). What is \emph{not} yet fixed
is which member of the PII family --- the value $\alpha(\beta)$ --- and hence the precise
$\beta$-dressing. That reduces O1$^\star$ from ``is the tail Painlev\'e~II?'' (now: yes) to ``which
$\alpha(\beta)$?'' (a bounded computation, Q1--Q2). This is a substantially smaller and sharper
statement than Step~1 began with, and it is the natural next write-up.

\begin{thebibliography}{9}
\bibitem{FN1980} H.~Flaschka, A.~C.~Newell, \emph{Monodromy- and spectrum-preserving deformations
I}, Comm. Math. Phys. \textbf{76} (1980) 65--116.
\bibitem{FIKN} A.~S.~Fokas, A.~R.~Its, A.~A.~Kapaev, V.~Yu.~Novokshenov, \emph{Painlev\'e
Transcendents: The Riemann--Hilbert Approach}, Math. Surveys Monogr. \textbf{128}, AMS, 2006.
\bibitem{Clarkson} P.~A.~Clarkson, \emph{Painlev\'e equations --- nonlinear special functions},
J. Comput. Appl. Math. \textbf{153} (2003) 127--140.
\bibitem{CleriDunne} N.~J.~Cleri, G.~V.~Dunne, \emph{Resurgent trans-series for generalized
Hastings--McLeod solutions}, J. Phys. A \textbf{53} (2020) 305201; arXiv:2002.06270.
\bibitem{Costin} O.~Costin, \emph{On Borel summation and Stokes phenomena of nonlinear
differential systems}, Duke Math. J. \textbf{93} (1998) 289--344; arXiv:math/0608408.
\bibitem{CLDS} A.~Comtet, P.~Le Doussal, N.~R.~Smith, \emph{The Tracy--Widom distribution at large
Dyson index}, arXiv:2510.14433 (2025).
\end{thebibliography}

\end{document}
